\chapter*{Перш ніж почати}

Ця книга --- про машину, яка важить півтора кілограма, споживає двадцять ватів, як тьмяна лампочка, і при цьому впізнає обличчя, вивчає мови й ловить м'яч на льоту. Про ваш мозок. Ми дивитимемося на нього так, як інженер дивиться на незнайому плату: які там деталі, як їх з'єднано, які сигнали біжать проводами і що ця схема обчислює.

Вам не знадобляться ні формули, ні знання біології. Знадобиться цікавість і готовність уявити собі провід, лампочку, вимикач. Усе інше ми зберемо дорогою.

Книга має другу, повну версію --- з рівняннями, моделями, задачами й посиланнями на досліди. Розділи обох версій ідуть в одному порядку й мають однакові номери. Якщо якийсь розділ вас зачепив, відкрийте його повну версію: там ви знайдете, звідки взялося кожне твердження і як його перевірити самому. На сайті книги кожен розділ можна перемикати між «Коротко» і «Повністю» однією кнопкою.

Одне застереження. Про мозок відомо далеко не все, і я чесно казатиму, де закінчується виміряне й починається здогад. Це не вада книги, а найцікавіше в ній: багато розділів закінчуються питаннями, на які поки ніхто не знає відповіді.
